\chapter*{\centering ВВЕДЕНИЕ}
\addcontentsline{toc}{chapter}{ВВЕДЕНИЕ}

\textbf{Актуальность.} 
В современной промышленности и материаловедении никелевые суперсплавы 
играют критически важную роль, особенно в авиастроении и энергетике, 
где требуются материалы, способные выдерживать экстремальные 
температуры и механические нагрузки. Традиционные методы подбора 
состава и прогнозирования свойств таких сплавов требуют проведения 
дорогостоящих и длительных физических экспериментов. В связи с этим, 
применение методов машинного обучения и искусственных нейронных сетей 
для ускорения процесса разработки новых материалов является 
высокоактуальной задачей.

\textbf{Проблема.} 
Несмотря на наличие больших массивов экспериментальных данных, 
классические эмпирические модели и физические уравнения часто не 
способны точно описать сложные нелинейные зависимости между 
химическим составом суперсплавов, параметрами термической обработки 
и их итоговыми физико-механическими свойствами. Это приводит к 
высокой погрешности прогнозов и необходимости многократного 
перебора вариантов при создании новых марок сплавов.

Целью работы является численное воспроизведение и анализ современных 
подходов к прогнозированию термических и механических свойств 
никелевых суперсплавов с использованием архитектур глубокого 
обучения и байесовских нейронных сетей.

Для достижения поставленной цели необходимо решить следующие задачи:
\begin{enumeratePaper}
    \item Изучить и воспроизвести модель глубокого обучения с 
    дифференциальным слоем для прогнозирования температуры сольвуса 
    никелевых суперсплавов.
    \item Реализовать байесовскую регуляризованную нейронную сеть (BRANN) 
    для прогнозирования предела прочности суперсплавов в зависимости 
    от параметра Ларсона-Миллера. 
    \item Провести сравнительный анализ точности полученных нейросетевых 
    моделей с оригинальными данными из научных статей.
\end{enumeratePaper}

\textbf{Краткое содержание и результаты.} 
В ходе выполнения работы был предобработан датасет, описывающий 
химический состав и свойства никелевых суперсплавов. Были реализованы две нейросетевые 
архитектуры. В результате воспроизведения первой модели удалось достичь 
высокой точности предсказания температуры сольвуса.
Во второй части работы с помощью BRANN удалось частино аппроксимирована 
зависимость предела прочности от параметра Ларсона-Миллера, но 
оптимальные параметры не были подобраны. Полученные результаты 
подтверждают эффективность применения методов машинного обучения в 
физике металлов.
\newpage